\documentclass[portrait,a0b,final]{a0poster}
\usepackage{pstricks,pst-grad}
\usepackage{graphicx}
\usepackage{url}
\usepackage{csquotes}
\MakeOuterQuote{�}
\usepackage{setspace}
% \usepackage{showframe}

% define poster column
\newenvironment{postercolumn}[1]{
   \begin{minipage}[t]{#1\textwidth}
   \begin{center}
}{
   \end{center}
   \end{minipage}
}

% define command for boxes
\newcommand{\posterbox}[4]{
\psshadowbox[#3]{
\begin{minipage}[t][#2][t]{#1}
\vspace{10mm}
\begin{spacing}{1.2}
#4
\end{spacing}
\end{minipage}
}}

% define some colors
\definecolor{lightblue}{rgb}{0, 0, 0.8}
\definecolor{white}{rgb}{1, 1, 1}
\definecolor{whiteblue}{rgb}{0.75, 0.8, 1}

% define font sizes
\newcommand{\standardfontsize}{\LARGE}
\newcommand{\boxheaderfontsize}{\huge}

% define font family
\renewcommand{\familydefault}{\sfdefault}

\setlength{\fboxsep}{0pt}


\begin{document}
% set background color (gradient)
\psframe[fillstyle=gradient, gradbegin=white, gradend=whiteblue,
   gradmidpoint=1.0](-3,3)(1.1\textwidth,-1.1\textheight)

\begin{center}
% \fbox{
\begin{postercolumn}{1}
% main header
\posterbox{0.96\textwidth}{}{linewidth=2mm, framearc=0.3,
   linecolor=lightblue, fillstyle=gradient, gradangle=0,
   gradbegin=white, gradend=whiteblue, gradmidpoint=1.0,
   framesep=1em}{
% left useR logo
\begin{minipage}[c][9cm][c]{0.2\textwidth}
  \begin{center}
    \includegraphics[width=12cm]{useR-large}
  \end{center}
\end{minipage}
% title etc
\begin{minipage}[c][13cm][c]{0.59\textwidth}
  \begin{center}
    {\VeryHuge \textbf{micEcon}}\\[12mm]
    {\Huge Microeconomic Analysis with R}\\[8mm]
    {\Huge Arne Henningsen and Ott Toomet}\\[8mm]
    {\LARGE University of Kiel (Germany) and University of Aarhus (Denmark)}
  \end{center}
\end{minipage}
% right useR logo
\begin{minipage}[c][9cm][c]{0.2\textwidth}
  \begin{center}
    \includegraphics[width=12cm]{useR-large}
  \end{center}
\end{minipage}
}
\end{postercolumn}

\vspace{15mm}

% \fbox{
\begin{postercolumn}{0.49}
% Demand Analysis
\posterbox{0.92\textwidth}{}{linewidth=2mm, framearc=0.1,
   linecolor=lightblue, fillstyle=gradient, gradangle=0,
   gradbegin=white, gradend=whiteblue, gradmidpoint=1.0,
   framesep=1em}{
{\boxheaderfontsize \textbf{Demand Analysis with the
�Almost Ideal Demand System� (AIDS)}}\\[5mm]
{\standardfontsize
The AIDS is the most popular demand system in empirical demand analysis.
\texttt{micEcon} provides functions for:
\medskip
\begin{itemize}
\item econometric estimation using the
\begin{itemize}
\item popular but inaccurate �Linear Approximation� (LA-AIDS)
\item accurate �Iterated Linear Least Squares Estimator� (ILLE)
\end{itemize}
\item searching for the intercept of the translog price index ($\alpha_0$)
   that gives the best fit to the model
\item optional imposition homogeneity and symmetry
\item checking for monotonicity and concavity
\item predicted quantities and expenditure shares
\item Marshallian (uncompensated) price elasticities, 
   Hicksian (compensated) price elasticities, 
   income elasticities,
   and their standard errors
\end{itemize}
\vspace{8mm}
}}

\vspace{15mm}

% SNQ Profit Function
\posterbox{0.92\textwidth}{}{linewidth=2mm, framearc=0.1,
   linecolor=lightblue, fillstyle=gradient, gradangle=0,
   gradbegin=white, gradend=whiteblue, gradmidpoint=1.0,
   framesep=1em}{
{\boxheaderfontsize \textbf{Production Analysis with the
�Symmetric Normalized Quadratic� (SNQ) profit function}}\\[5mm]
{\standardfontsize
The SNQ profit function is often used in microeconomic production models
because it enables the imposition of global convexity.
\texttt{micEcon} provides functions for:
\begin{itemize}
\item econometric estimation
\item imposition of global convexity using a recent two-step estimation
   procedure
\item price elasticities of inputs and outputs and their standard errors
\item fixed factor elasticities
\item shadow prices of fixed factors
\item Hessian matrix and
   its derivatives with respect to coefficients
\item predicted profits and netput quantities
   including the standard errors and the confidence limits of prediction
\end{itemize}
\vspace{8mm}
}}
\end{postercolumn}
% }
\hfill
% \fbox{
\begin{postercolumn}{0.49}
% Sample selection models
\posterbox{0.92\textwidth}{}{linewidth=2mm, framearc=0.1,
   linecolor=lightblue, fillstyle=gradient, gradangle=0,
   gradbegin=white, gradend=whiteblue, gradmidpoint=1.0,
   framesep=1em}{
{\boxheaderfontsize \textbf{Heckman Selection Models}}\\[5mm]
{\standardfontsize
If a non-random (sub)sample is used for econometric estimation,
Heckman selection models are required to obtain unbiased estimates.
\texttt{micEcon} provides functions for:
\medskip
\begin{itemize}
\item two-step Heckman estimation
\item Maximum Likelihood estimation
\item variance-covariance matrices for
Heckman-type estimations
\item �Inverse Mill's Ratios� for standard selection models
(from univariate probit estimations)
\item �Inverse Mill's Ratios� for double-selection models
(from bivariate probit estimations)
\end{itemize}
}}

\vspace{15mm}

% Matrix Tools
\posterbox{0.92\textwidth}{}{linewidth=2mm, framearc=0.1,
   linecolor=lightblue, fillstyle=gradient, gradangle=0,
   gradbegin=white, gradend=whiteblue, gradmidpoint=1.0,
   framesep=1em}{
{\boxheaderfontsize \textbf{Tools for Matrices}}\\[5mm]
{\standardfontsize
\texttt{micEcon} provides functions for:
\medskip
\begin{itemize}
\item inserting additional rows and columns into matrices
\item creating symmetric and triangular matrices
\item conversion between a symmetric matrix and a vector
   of its linear independent values
\item tests for positive and negative semidefiniteness
% \item checking quasiconcavity and quasiconvexity
\end{itemize}
}}

\vspace{15mm}

% Miscellaneous Tools
\posterbox{0.92\textwidth}{}{linewidth=2mm, framearc=0.1,
   linecolor=lightblue, fillstyle=gradient, gradangle=0,
   gradbegin=white, gradend=whiteblue, gradmidpoint=1.0,
   framesep=1em}{
{\boxheaderfontsize \textbf{Miscellaneous Tools}}\\[5mm]
{\standardfontsize
\texttt{micEcon} provides functions for:
\medskip
\begin{itemize}
\item quadratic functions: e.g.\ estimation,
   derivatives with respect to regressors
\item translog functions: e.g.\ estimation,
   derivatives and (bordered) Hessian matrix with respect to regressors
\item write input files and read output files of
FRONTIER~4.1 (software for stochastic frontier analysis)
\item maximization using the Newton-Raphson and BHHH algorithm
\end{itemize}
}}

\end{postercolumn}
% }

\vspace{15mm}

\begin{postercolumn}{1}
\posterbox{0.96\textwidth}{}{linewidth=2mm, framearc=0.3,
   linecolor=lightblue, fillstyle=gradient, gradangle=0,
   gradbegin=white, gradend=whiteblue, gradmidpoint=1.0,
   framesep=1em}{
\begin{spacing}{1.5}
{\Large
\textbf{Arne Henningsen} $\cdot$
Department of Agricultural Economics $\cdot$
University of Kiel $\cdot$
Olshausenstr.~40 $\cdot$
D-24098 Kiel (Germany) $\cdot$
Phone~+49-431-880 4445 $\cdot$
Fax~+49-431-880 1397 $\cdot$
\url{ahenningsen@agric-econ.uni-kiel.de} $\cdot$
\url{http://www.uni-kiel.de/agrarpol/ahenningsen/}\\[5mm]
\textbf{Ott Toomet} $\cdot$
Department of Economics $\cdot$
University of Aarhus $\cdot$
Building~1322 $\cdot$
DK-8000 \AA rhus C (Denmark) $\cdot$
Phone~+45-8942-1589 $\cdot$
Fax~+45-8613-6334 $\cdot$
\url{otoomet@econ.au.dk} $\cdot$
\url{http://www.econ.au.dk/afn/PhD/toometo.htm}%
}%
\end{spacing}%
}
\end{postercolumn}
% }
\end{center}
\end{document}
